\chapter{Alcances del sistema}
\label{cap:alcances}

El proyecto comprenderá el análisis del proceso de venta y reserva de
paquetes turísticos de Vizcacha Tours y el desarrollo de un sistema web que
permita administrar paquetes, programar salidas, consultar la disponibilidad
de cupos, registrar clientes y pasajeros, elaborar cotizaciones, gestionar
reservas, registrar anticipos y saldos, administrar proveedores y generar
reportes. La prioridad será verificar los cupos y registrar correctamente las
reservas; las demás funciones se incorporarán como apoyo a ese proceso
principal.

\section{Módulos propuestos}

\begin{itemize}
  \item \textbf{Paquetes turísticos:} registro de destinos, descripción,
        duración, itinerario, precios y servicios incluidos.
  \item \textbf{Salidas y disponibilidad:} programación de fechas, capacidad
        de pasajeros y cálculo de cupos disponibles.
  \item \textbf{Clientes y pasajeros:} registro de datos necesarios para
        atender y dar seguimiento a las solicitudes.
  \item \textbf{Cotizaciones:} preparación y consulta de propuestas según el
        paquete solicitado y sus condiciones.
  \item \textbf{Reservas:} creación, consulta, confirmación, modificación y
        cancelación con validación de cupos.
  \item \textbf{Pagos:} registro de anticipos y saldos, con pago en efectivo y
        simulación de banca móvil o código QR, sin integración con pasarelas
        bancarias reales.
  \item \textbf{Proveedores:} datos básicos de transporte, hospedaje y guías
        vinculados a los paquetes.
  \item \textbf{Reportes y usuarios:} consultas por fecha, ocupación, reservas
        y accesos según el rol de cada trabajador.
\end{itemize}

Para facilitar su uso, se procurará una interfaz sencilla y adaptable a
distintos tamaños de pantalla.
